\MMHostGuidanceHeading


\subsubsection{Being the game master}

As the GM, you should balance two priorities:
\begin{itemize}
  \item The game should be fun.
  \item The game should be fair.
\end{itemize}

Give players information at the right time. Context their character would have known for a long time can be shared in advance, ideally at least a day before the session. Deliver recent, sensitive, or plot-critical information later, when it becomes relevant to play.

\subsubsection{The flow of the game}

Try to separate the game in phases:
\begin{enumerate}
  \item Introductions, in which everyone establishes their character.
  \item A first period of social play, relationships, and low-stakes hints.
  \item An inciting event that makes the main mystery unmistakable.
  \item Investigation, private consultations with the GM when needed, and a final vote or accusation.
\end{enumerate}


\subsubsection{Writing the story}

Start with the answer to the main mystery: who committed the act, what happened, or what the supernatural or hidden truth really is. Build the story backwards from that answer. Connect every character to at least one other character, and give at least half of them information that can contribute to solving the main mystery.

To write the story, start by writing the atomic facts~\ref{subsec:atomic_facts}. These will be the facts from an outside perspective. Keep in mind that players may lie during the game, or never find some of the facts. Because of this, make it so that, if a player had access to all these facts, finding the truth would be trivial.

Then, write the background~\ref{subsec:background}. This is also the outside perspective, and is mainly the more well-written version of all the facts. As you write this, you may want to come back and change the facts, as you'll gain a new perspective on the story.

As you write these two portions on the game, be sure to create a complete version of the characters, then of the mysteries and then of the clues.

Only after you are confident on the current state of the story should you start writing the story from the perspective of each character. In the character's stories, you may want to embellish it more, show more of their inner thoughts, more of their personality, and you may also add some unnecessary information. You may still go back and change the whole story, but be sure to leave all the texts consistent.

\subsubsection{Creating the characters}

Regarding creating the characters, remember that, besides the gm notes, that text will be shown to the players. Most of the text they have should be with their POVs, embellished with their opinions and personality. There should always exist some extra unnecessary information.

Take notice on the points you award regarding their difficulty.

\subsubsection{Game balance}

Before thinking about balance, remember that the game should be fun. If the players are having a good time finding and solving clues while arguing, it is all good.

Regarding difficulty, don't try to make the mystery overly complex. You are not writing a book or directing a series of which the players will have countless hours to interpret. There will be a lot of distractions.

If someone had access to almost 100\% of the facts, the answer to the mystery should be obvious.
If someone had access to 75\% of the facts and half an hour to think, it should be fairly easy to resolve the underlying mystery. If someone had access to 50\% of the facts and an hour to think, the correct answer should at least be possible to reach.

Take into account time lost in red herrings, the time players will take trying to find out each other secrets, and so on.

As clues are found, the story the players will construct should naturally flow in a relevant way for them to solve the underlying mystery.
